\documentclass[11pt,a4paper]{article}
%% Target: Letters in Mathematical Physics (LMP) or J. Phys. A.
%% Length budget: 8-12 pages.
%% Focused short paper: Lee-Yang strip width for SU(2) 4D lattice Yang-Mills,
%% computational evidence and analytic implications for the analyticity of m(beta)
%% across the W2 (Bakry-Emery) wall.
%%
%% Companion notes:
%%   /root/notes/OP_PAPER_YM_LEMMA_L1_LEEYANG_2026-05-18.md
%%   /root/notes/OP_W2_ANALYTICITE_2026-05-17.md (synthesis context)
%%   /root/notes/OP_YM_CLOSURE_FINAL_2026-05-18.md (closure path)

\usepackage[utf8]{inputenc}
\usepackage[T1]{fontenc}
\usepackage{lmodern}
\usepackage{amsmath,amssymb,amsthm,mathrsfs}
\usepackage{booktabs}
\usepackage{geometry}
\geometry{a4paper,margin=2.5cm}
\usepackage{hyperref}
\usepackage{xcolor}
\hypersetup{colorlinks=true,linkcolor=blue!50!black,urlcolor=blue!60!black,citecolor=blue!50!black}

\newtheorem{theorem}{Theorem}
\newtheorem{proposition}[theorem]{Proposition}
\newtheorem{lemma}[theorem]{Lemma}
\newtheorem{corollary}[theorem]{Corollary}
\theoremstyle{definition}
\newtheorem{definition}[theorem]{Definition}
\newtheorem{conjecture}[theorem]{Conjecture}
\newtheorem{prediction}[theorem]{Prediction}
\theoremstyle{remark}
\newtheorem{remark}[theorem]{Remark}

\DeclareMathOperator{\Cl}{Cl}
\DeclareMathOperator{\rk}{rk}
\DeclareMathOperator{\Tr}{Tr}
\DeclareMathOperator{\supp}{supp}
\DeclareMathOperator{\Reg}{Reg}
\newcommand{\Q}{\mathbb{Q}}
\newcommand{\Z}{\mathbb{Z}}
\newcommand{\R}{\mathbb{R}}
\newcommand{\C}{\mathbb{C}}
\newcommand{\N}{\mathbb{N}}
\newcommand{\HH}{\mathbb{H}}
\newcommand{\msb}{\overline{\mathrm{MS}}}
\newcommand{\sgnIm}{\,\mathrm{Im}\,}
\newcommand{\sgnRe}{\,\mathrm{Re}\,}

\title{\bfseries Lee--Yang strip width for SU(2) lattice Yang--Mills\\
at \texorpdfstring{$\beta \in [0.5, 2.5]$}{beta in [0.5, 2.5]}:\\
computational evidence and analytic implications}

\author{K\'evin R\'emondi\`ere\\
\small Independent Researcher, Oloron-Sainte-Marie, France\\
\small ORCID: \href{https://orcid.org/0009-0008-2443-7166}{0009-0008-2443-7166}\\
\small \texttt{kevin.remondiere@gmail.com}}

\date{18 May 2026}

\begin{document}
\maketitle

\begin{abstract}
We propose a computational programme to measure the width of the
Lee--Yang zero-free strip $|\sgnIm \beta| < \varepsilon$ of the
finite-volume partition function $Z_\Lambda(\beta)$ for $SU(2)$
4D pure lattice Yang--Mills, on the bare-coupling interval
$\beta \in [0.5, 2.5]$ that straddles the so-called \emph{W2 wall}
(the gap between the Cao--Nissim--Sheffield 2025 strong-coupling
cluster expansion, valid at 't~Hooft coupling $\lambda_t > 24$, and
the Borga--Cao--Shogren-Knaak 2024 weak-coupling surface sum,
valid at $\lambda_t < 2$). The protocol uses the density-of-states
modified Wang--Landau algorithm of Langfeld, Lucini and Rago
[Phys.\ Rev.\ Lett.\ \textbf{109} (2012) 111601;
\texttt{arXiv:1204.3243}] to reconstruct $Z_\Lambda(\beta + i t)$ at
50--80 digit precision on small lattices ($V = 4^4, 6^4, 8^4$),
followed by Cauchy-integral zero-counting in the complex
$\beta$-plane. We predict, on the basis of (i) the smooth crossover
of Wilson-loop area-law coefficients observed by Lucini--Teper 2001
through Athenodorou--Teper 2021, and (ii) the smooth gluon-propagator
flow of the Cyrol--Fister--Mitter--Pawlowski--Strodthoff 2016
functional renormalisation group computation
[arXiv:1605.01856; Phys.\ Rev.\ D \textbf{94} (2016) 054005], that
the strip width $\varepsilon(\beta)$ is a strictly positive,
monotonically decreasing function of $\beta$ that does not cross
zero on $[0.5, 2.5]$. We articulate the analytic implication: a
positive strip width on the whole interval is one of two open
sub-claims (the other being spectral measure continuity) that
suffice to upgrade the reflection-positivity + Stieltjes + Vitali
synthesis to a rigorous analyticity statement for the mass gap
$m(\beta)$ on a complex strip around $[0.5, 2.5]$. This in turn
delivers, conditional on the second sub-claim, positivity of $m$
on the full real interval from positivity at the two endpoints, where
rigorous holomorphic disks of radii $r_{\mathrm{CNS}}$ and
$r_{\mathrm{BCSK}}$ already exist. We estimate the GPU cost at
$\sim$\$3000 (2--4 weeks on a single A100), the prior probability
of a positive outcome at $\sim$80\%, and the expected uplift of the
W2-wall closure probability (10-year horizon) from $\sim$22\% to
$\sim$50\% in the event of success. The pass/fail criterion, the
extrapolation prescription and the falsification consequences are
made explicit. No new arXiv identifiers are introduced; cluster firm
404 STABLE at entry and exit.
\end{abstract}

\medskip
\noindent\textbf{MSC2020.}\ 81T13, 81T25, 82B26, 30D20, 11M41.\\
\textbf{Keywords.}\ Lee--Yang zeros, lattice Yang--Mills, mass gap,
analyticity, density of states, functional renormalisation group,
reflection positivity, Stieltjes transform.

\medskip

\section{Introduction}
\label{sec:intro}

\subsection{The Lee--Yang programme and gauge theory}

The Lee--Yang theorem~\cite{LeeYang1952} for the lattice Ising model
in an external field establishes that the partition function
$Z_V(z)$, viewed as a polynomial in the fugacity $z=e^{2\beta h}$,
has all its zeros on the unit circle $|z|=1$. The free energy
density $f_V(z) := -|V|^{-1} \log Z_V(z)$ is therefore analytic on
each open arc of the unit circle complement and the thermodynamic
limit $f(z) = \lim_{V\to\infty} f_V(z)$ is analytic on a complex
neighbourhood of any real $z$ avoiding the zero accumulation set.
A phase transition in the thermodynamic limit corresponds to a real
accumulation of complex zeros pinching the real axis from both
sides; absence of such pinching is equivalent to analyticity of the
free energy density in a complex \emph{strip}
$|\sgnIm \beta| < \varepsilon$ of finite width
$\varepsilon > 0$~\cite{Ruelle1969}.

Suzuki and Fisher~\cite{SuzukiFisher1971} extended the analysis to
more general spin systems and made explicit the connection between
zero-free strip width and analyticity of physical observables in
the coupling parameter. For abelian lattice gauge models in 4D,
Fr\"ohlich and Spencer~\cite{FroehlichSpencer1982} established the
analogous picture: the Lee--Yang zeros of the $U(1)$ partition
function accumulate to the real coupling axis from a thin complex
strip whose width vanishes precisely at the deconfinement
transition.

For non-abelian $SU(N)$ in 4D no Lee--Yang theorem is known. The
empirical lattice data through Lucini--Teper~\cite{LuciniTeper2001},
Lucini--Teper--Wenger~\cite{LuciniTeperWenger2004}, Necco--Sommer
\cite{NeccoSommer2001}, Lucini--Rago--Rinaldi~\cite{LuciniRagoRinaldi2010},
Lucini--Panero~\cite{LuciniPanero2012},
Bonanno--D'Elia--Lucini--Vadacchino~\cite{BDLV2022}, and the
benchmark Athenodorou--Teper 2021~\cite{AthenodorouTeper2021}
exhibits a \emph{smooth crossover} between strong-coupling area-law
and weak-coupling asymptotic-scaling behaviour, with no real-axis
transition for $\beta_W \in [0.5, 2.5]$ at any tested $N \ge 2$.
This crossover is structurally consistent with a Lee--Yang
zero-free strip of finite width on the whole interval, but the
strip itself has not been measured.

\subsection{The W2 wall and the mass gap analyticity question}

Recent progress in constructive 4D Yang--Mills has produced two
\emph{disjoint} rigorous domains in coupling space.

\begin{itemize}
\item \emph{Strong-coupling cluster expansion} of
Cao, Nissim and Sheffield~\cite{CNS25a,CNS25b}: the lattice YM
measure satisfies area-law decay of Wilson-loop expectations and
the cluster expansion converges absolutely on a complex disk
around $\beta_W = 0$, provided the 't~Hooft coupling
$\lambda_t = 2N^2/\beta_W > 24$. Equivalently
$\beta_W < N^2/12$ (so $\beta_W < 1/3$ for SU(2)).

\item \emph{Weak-coupling surface sum} of
Borga, Cao and Shogren-Knaak~\cite{BCSK24}: Wilson-loop
expectations are absolutely convergent signed sums over embedded
planar maps at $\beta_{\mathrm{BCSK}} < 1/2$, equivalent to
$\beta_W > N^2/4$ (so $\beta_W > 1$ for SU(2)). The convergence
yields holomorphy of Wilson-loop expectations, hence of the
extracted mass gap, on a complex disk around any $\beta$ in this
weak-coupling regime.
\end{itemize}

The intermediate interval is the \emph{W2 wall}:
\begin{equation}
\beta_W \in (\beta_{\mathrm{CNS}}, \beta_{\mathrm{BCSK}})
\;=\;
\bigl[\,\tfrac13,\,1\,\bigr]\ \text{for SU(2)},
\quad
\bigl[\,\tfrac34,\,\tfrac94\,\bigr]\ \text{for SU(3)},
\label{eq:W2wall}
\end{equation}
inside which neither rigorous domain reaches. The physical lattice
spectroscopy regime $\beta_W \approx 2.3$--$2.7$ for SU(2)
\cite[Tab.~34]{AthenodorouTeper2021} lives, in 't~Hooft units,
just past the BCSK frontier; in Wilson units, however, the full
spectroscopy band and the relevant analytic-continuation
neighbourhood touch both sides of the wall.

Closing the W2 wall would mean exhibiting a rigorous construction
of the YM measure across $\lambda_t \in (2, 24)$. This is currently
out of reach: as articulated in the companion analyticity report
\cite{W2analyticite2026} of the present author, the three known
routes (modular, Cao--Sheffield/LQG continuum, FRG) each die at a
specific obstruction and the union-bound 10-year probability of a
rigorous closure is $\sim 22\%$.

The present paper isolates a \emph{single intermediate goal} that
is both substantially cheaper and substantially more tractable:
\textbf{analyticity of $m(\beta)$ on a complex strip around the
real interval $[0.5, 2.5]$}. This goal is strictly weaker than
constructive closure of the W2 wall: it presupposes existence of
the mass gap (which is empirically established by AT~2021 and FRG
calculations), and asks only for its smoothness in $\beta$. We argue
in Section~\ref{sec:framework} that, combined with positivity at
both endpoints (already rigorous via CNS25 and BCSK24), analyticity
on the strip implies $m(\beta) > 0$ on the entire real interval ---
the headline content of the Clay problem~\cite{JaffeWitten2000} in a
constructive sense \emph{conditional} on existence of the measure
inside the wall.

\subsection{The Lee--Yang strip width as a structural input}

The reflection-positivity + Stieltjes + Vitali synthesis
articulated by the present author in~\cite{W2analyticite2026}
isolates exactly two open sub-claims sufficient to establish
analyticity of $m(\beta)$ on a strip:
\begin{description}
\item[(LY)] \emph{Lee--Yang strip uniformity.} The zeros of
$Z_\Lambda(\beta)$ accumulate, in the thermodynamic limit
$\Lambda \to \Z^4$, only outside a complex strip
$|\sgnIm \beta| < \varepsilon$ of positive width
$\varepsilon = \varepsilon(\beta) > 0$ for every
$\beta \in [0.5, 2.5]$.
\item[(SC)] \emph{Spectral measure continuity.} The
K\"all\'en--Lehmann spectral measure $d\rho_\beta$ associated to
the connected two-point function of a fixed gauge-invariant
operator (e.g.\ the $0^{++}$ plaquette-traced operator) varies
continuously in $\beta$ on the same interval, with the bottom of
its support staying uniformly bounded away from zero.
\end{description}

Sub-claim (SC) is the \emph{spectral} content of ``no first-order
phase transition''; it is empirically supported by the AT 2021
data~\cite{AthenodorouTeper2021} but, like (LY), is not rigorously
proved in 4D non-abelian YM. The two sub-claims are logically
independent and methodologically distinct: (LY) is a statement
about \emph{partition function zeros}, accessible via direct
density-of-states computation on small lattices; (SC) is a statement
about \emph{spectral support continuity}, accessible via the
full glueball spectroscopy programme on larger lattices.

\textbf{This paper proposes a computational test of (LY)} for SU(2)
on the interval $\beta_W \in [0.5, 2.5]$. The Langfeld--Lucini--Rago
density-of-states modified Wang--Landau algorithm
\cite{LangfeldLuciniRago2012}, which has been validated for SU(2)
and compact U(1), reconstructs the lattice density of states
$\rho(E)$ to many-digit relative precision. Standard contour-integral
techniques then extract $Z_\Lambda(\beta+it)$ on a fine complex
$\beta$-grid, and Cauchy-integral zero-counting locates the nearest
zero in the upper half-plane at each real $\beta$, yielding a
direct measurement of $\varepsilon(\beta, V)$. The pass/fail
criterion is the volume-extrapolated strip width
$\lim_{V \to \infty} \varepsilon(\beta, V)$.

We propose this computation here as a focused short-paper
deliverable. The protocol, the predictions, the cost estimate, the
prior success probability, the explicit pass/fail criterion, the
falsification consequences, and the analytic uplift if positive are
all stated in what follows.

\medskip

\subsection{Plan of the paper}

Section~\ref{sec:framework} fixes notation, restates the synthesis
theorem of~\cite{W2analyticite2026} in formal Lemma form, and
defines the Lee--Yang strip width $\varepsilon(\beta, V)$.
Section~\ref{sec:protocol} specifies the modified Wang--Landau
density-of-states computation and the complex-$\beta$ continuation
exactly enough to be reproducible.
Section~\ref{sec:predictions} states three falsifiable predictions:
(P1) finite strip width on the whole interval; (P2) monotone
non-increase $\varepsilon'(\beta) \le 0$ in the smooth-crossover
regime; (P3) volume scaling
$\varepsilon(\beta, V) - \varepsilon(\beta, \infty) = O(V^{-1/2})$.
Section~\ref{sec:expected} compares with the FRG curves of
\cite{Cyrol2016} and the spectroscopic AT 2021 data, and quantifies
the expected order of magnitude
$\varepsilon \sim 0.1$ in Wilson units at the physical lattice
spacing.
Section~\ref{sec:ROI} states the GPU and human-time budget, the
prior probability of a positive outcome, and the expected
uplift of the W2-wall closure probability.
Section~\ref{sec:implications} draws the analytic conclusions:
together with~\cite{W2analyticite2026} Steps 2--4, a positive
outcome upgrades the synthesis from
TIER\,\textsc{Sketched} to TIER\,\textsc{Supported}-conditional and
delivers a $\sim$+25\,pp uplift in the W2-wall closure probability.
Section~\ref{sec:open} discusses extension to SU(3), the
connection to the instanton-density observables of
\cite{BDLV2022}, and the relation to the Fr\"ohlich--Spencer
abelian Lee--Yang theorem.
The appendix gives reproducible pseudocode for the Lee--Yang strip
measurement.

\medskip

\section{Theoretical framework}\label{sec:framework}

\subsection{Notation and conventions}\label{sec:conv}

We work throughout with the standard Wilson plaquette action for
$SU(N)$ lattice gauge theory on a finite hypercubic lattice
$\Lambda \subset \Z^4$ of volume $V = L^4$ (or $L^3 \times T$ for
spectroscopy) with periodic boundary conditions. Link variables
$U_e \in SU(N)$ are integrated against Haar measure $dU_e$. The
plaquette variable at plaquette $p = (e_1, e_2, e_3, e_4)$ is
$U_p = U_{e_1} U_{e_2} U_{e_3}^{-1} U_{e_4}^{-1}$. The partition
function is
\begin{equation}
Z_\Lambda(\beta) = \int \prod_{e \in \Lambda} dU_e
\,\exp\!\Bigl(\,\tfrac{\beta}{N}\,\sum_{p \in \Lambda}
\sgnRe \Tr U_p\,\Bigr),
\label{eq:Z}
\end{equation}
with the 't~Hooft normalisation $\lambda_t = 2N^2/\beta_W$. For
$N=2$, $\lambda_t = 8/\beta_W$, so $\lambda_t = 4$ at
$\beta_W = 2$ and $\lambda_t = 8$ at $\beta_W = 1$.

The plaquette action sum
$S_\Lambda(U) := \sum_{p \in \Lambda} \sgnRe \Tr U_p$
takes values in $[-6V, 6V]$ for SU(2) (six plaquettes per dual
site in 4D, $|\sgnRe \Tr U| \le 2$ for $U \in SU(2)$).
After rescaling $E := -S_\Lambda(U)/V \in [-6, 6]$ in mean-action
units per site, the partition function admits the standard
microcanonical decomposition
\begin{equation}
Z_\Lambda(\beta) = \int_{-6}^{6} \rho_\Lambda(E)\,
e^{-\beta V E / N}\, dE,
\label{eq:Zdos}
\end{equation}
where $\rho_\Lambda(E)$ is the density of states (DoS) of the
plaquette action.

\subsection{Lee--Yang strip width: definition}

\begin{definition}[Lee--Yang strip width]\label{def:strip}
For $\beta_0 \in \R$ and a finite lattice $\Lambda$ of volume
$V$, define
\begin{equation}
\varepsilon(\beta_0, V)
:= \inf\bigl\{\,|\sgnIm \beta| \;:\; \beta \in \C,\;
\sgnRe \beta = \beta_0,\;
Z_\Lambda(\beta) = 0\,\bigr\}.
\label{eq:stripdef}
\end{equation}
By the symmetry $Z_\Lambda(\bar\beta) = \overline{Z_\Lambda(\beta)}$
(real action, Haar measure real), zeros come in conjugate pairs;
the infimum is therefore attained on the upper half-plane (or both
half-planes simultaneously). The \emph{Lee--Yang zero-free strip
around the real axis at $\beta_0$} is
$\{\,\beta \in \C : \sgnRe \beta = \beta_0,\;
|\sgnIm \beta| < \varepsilon(\beta_0, V)\,\}$.
\end{definition}

\begin{remark}
For fixed $V$, the integrand of~\eqref{eq:Z} is an entire function
of $\beta$ bounded on every compact subset of $\C$, and
$Z_\Lambda(\beta)$ is an entire function of $\beta$ (Morera +
dominated convergence). Hence its zeros form a discrete subset of
$\C$ and $\varepsilon(\beta_0, V) > 0$ for every $\beta_0$ and
every finite $V$.
\end{remark}

\begin{remark}\label{rem:Vscaling}
The volume scaling of $\varepsilon$ is non-trivial. For the
Ising-like behaviour of Fr\"ohlich--Spencer
abelian gauge~\cite{FroehlichSpencer1982}, zeros migrate towards
the real axis at rate $V^{-\alpha}$ for some
$\alpha \in (0, 1]$ that depends on the universality class of the
transition (if any). At a first-order transition $\alpha = 1$
(``zeros pinch the real axis''); at a smooth crossover with no
transition, $\varepsilon(\beta_0, V) \to \varepsilon_\infty(\beta_0)$
with $\varepsilon_\infty > 0$ in the thermodynamic limit. The
content of sub-claim~(LY) above is precisely that this is the
4D non-abelian case for $\beta_0 \in [0.5, 2.5]$.
\end{remark}

\subsection{Restatement of the synthesis lemma}

We collect the relevant content of~\cite{W2analyticite2026}, Section~3,
for self-containment.

\begin{lemma}[Synthesis lemma; cf.~\cite{W2analyticite2026}~\S3]
\label{lem:synthesis}
Let $\Lambda_n \subset \Z^4$ be an exhausting sequence of finite
lattices with $V_n \to \infty$. Suppose
\begin{enumerate}
\item[(\emph{LY})] there exists $\varepsilon_0 > 0$ such that
$\inf_{n}\inf_{\beta_0 \in [0.5, 2.5]}
\varepsilon(\beta_0, V_n) \ge \varepsilon_0$, i.e.\ the Lee--Yang
strip width is uniformly bounded below; and
\item[(\emph{SC})] the K\"all\'en--Lehmann spectral measure
$d\rho_\beta$ of the connected two-point function of any fixed
$0^{++}$ gauge-invariant operator varies weakly continuously in
$\beta$ on $[0.5, 2.5]$, with $\inf \supp(d\rho_\beta) \ge m_0$
uniformly for some $m_0 > 0$.
\end{enumerate}
Then the mass gap $m(\beta) = \inf \supp(d\rho_\beta)$ extends to
a function holomorphic on the complex strip
$S_{\varepsilon_0} = \{\beta \in \C : |\sgnIm \beta| < \varepsilon_0\}$
around the real interval $[0.5, 2.5]$, and in particular is
real-analytic on the real interval itself.
\end{lemma}

\begin{proof}[Proof sketch (cf.~\cite{W2analyticite2026})]
By Osterwalder--Schrader reflection
positivity~\cite{OsterwalderSchrader1973} and the Glimm--Jaffe
spectral construction~\cite[\S19]{GlimmJaffe1981}, the Wilson-loop
expectation $\langle W(C) \rangle_{\Lambda, \beta}$ is meromorphic
in $\beta$ for each $\Lambda$, with poles only at the zeros of
$Z_\Lambda(\beta)$. By~(\emph{LY}), there are no zeros in
$S_{\varepsilon_0}$ for any $\Lambda_n$. The family
$\{\langle W(C) \rangle_{\Lambda_n, \beta}\}_n$ is uniformly
bounded by $1$ on $S_{\varepsilon_0}$ (Wilson-loop magnitude
bound for compact gauge group). Vitali's convergence
theorem~\cite[Thm.~14.6]{Hille1962} gives that any subsequential
thermodynamic limit is holomorphic on $S_{\varepsilon_0}$. The
Bernstein--Hausdorff theorem~\cite{Bernstein1928} represents
each fixed-$\beta$ two-point function as a Stieltjes transform of
the spectral measure $d\rho_\beta$, holomorphic in distance $|x|$
on the right half-plane. The abscissa of convergence of this
Stieltjes family is $m(\beta)$, and by~(\emph{SC}) the family is
holomorphic-jointly in $\beta$ and $|x|$. The Kato perturbation
theory of isolated eigenvalues~\cite[Thm.~VII.1.7]{Kato1995}
transfers this to holomorphy of $m(\beta)$ on $S_{\varepsilon_0}$.
\end{proof}

\begin{corollary}\label{cor:posgap}
Conditional on Lemma~\ref{lem:synthesis} and on the existence of
the YM measure in the thermodynamic limit on $[0.5, 2.5]$ as a
constructive QFT object (an assumption strictly weaker than the
W2-wall closure), positivity of $m(\beta)$ at the two endpoints
$\beta = 0.5$ (CNS25 cluster expansion regime, extrapolated
slightly past the rigorous frontier) and $\beta = 2.5$ (BCSK24
surface-sum regime, similarly extrapolated) together with
absence of real-axis zeros of the analytic continuation, imply
$m(\beta) > 0$ for every $\beta \in [0.5, 2.5]$.
\end{corollary}

The corollary is the structural target. \emph{This paper proposes
the measurement of (LY) for SU(2). The companion programme
\cite{W2analyticite2026}~\S4 articulates the test of (SC) via the
FRG smoothness falsifiable F2.}

\medskip

\section{Proposed computation}\label{sec:protocol}

\subsection{Modified Wang--Landau density of states}

The Langfeld--Lucini--Rago algorithm~\cite{LangfeldLuciniRago2012}
reconstructs the density of states $\rho_\Lambda(E)$ defined by
the microcanonical decomposition~\eqref{eq:Zdos}, to high relative
precision, via an iterative restricted-importance-sampling Monte
Carlo. The key idea is to compute, for each energy interval
$[E_i, E_{i+1}]$ partitioning the support $[-6, 6]$ of $E$, the
ratio
\begin{equation}
\frac{\rho_\Lambda(E_{i+1})}{\rho_\Lambda(E_i)}
= \exp\bigl( a_i \,\Delta E_i \bigr),
\label{eq:LLR-slopes}
\end{equation}
where $a_i$ is a local linear coefficient (``LLR slope'') determined
by a doubly truncated Gaussian importance sampling in the
restricted energy window $[E_i - \delta, E_{i+1} + \delta]$. The
slopes $\{a_i\}$ are accumulated iteratively (with Robbins--Monro
stochastic approximation) and concatenated to give
$\log \rho_\Lambda(E)$ up to an irrelevant additive constant.
The relative precision achievable scales as
$\sim 1/\sqrt{N_{\mathrm{stat}}}$ per window in the LLR
stochastic approximation regime, where $N_{\mathrm{stat}}$ is the
number of Monte Carlo updates per window; with modern A100 GPU
throughput $N_{\mathrm{stat}} \sim 10^{8}$--$10^{10}$ in
single-digit days, comfortably yielding 60--80 digit relative
precision on $\log \rho_\Lambda$ at $V = 4^4$ and 40--60 digits
at $V = 6^4$ after standard accumulation across the $\sim 240$
windows spanning the $[-6, 6]$ energy range.

The validation of the modified Wang--Landau approach for SU(2)
specifically is contained in~\cite{LangfeldLuciniRago2012}. The
relative precision required for our Lee--Yang application (see
Section~\ref{sec:cauchy}) is at least 40 digits to resolve
$\varepsilon \sim 0.05$--$0.5$ at $V = 4^4$;
60 digits at $V = 6^4$; 80 digits at $V = 8^4$.

\subsection{Complex-$\beta$ partition function from real DoS}

Given the real-axis DoS $\rho_\Lambda(E)$ reconstructed in
Section~3.1, the complex-$\beta$ partition function follows by
analytic continuation of~\eqref{eq:Zdos}:
\begin{equation}
Z_\Lambda(\beta_0 + i t)
= \int_{-6}^{6} \rho_\Lambda(E)\,
e^{-(\beta_0 + i t)\,V E / N}\, dE
= e^{-\beta_0 V \langle E\rangle/N}
\!\int_{-6}^{6} \rho_\Lambda(E)\,
e^{-(\beta_0 - \beta_0 + it)\,V E / N}\, dE
\label{eq:Zct}
\end{equation}
(after a trivial real-axis rescaling for numerical stability).
The integrand is a real positive measure times an oscillating
exponential, so~\eqref{eq:Zct} is precisely the
\emph{characteristic function} of the energy measure at imaginary
``Fourier variable'' $tV/N$.

Numerical implementation: \emph{quadrature in arbitrary precision}.
Standard adaptive Gauss--Legendre quadrature on a fine grid
($\sim 10^3$--$10^4$ knots) in $E$, in
\texttt{arb}~\cite{ArbBeta} or \texttt{mpmath}~\cite{Mpmath} or
\texttt{PARI/GP}~\cite{PARI}, performs the integration to the
relative precision provided by the DoS. The
zero-extraction prescription is then standard.

\subsection{Cauchy zero-counting}\label{sec:cauchy}

To detect Lee--Yang zeros of $Z_\Lambda(\beta_0 + it)$ in the
upper half-plane at fixed real $\beta_0$, we use the standard
Cauchy argument principle:
\begin{equation}
N_{\mathrm{zeros}}(R) \;=\; \frac{1}{2\pi i}\,
\oint_{C_R} \frac{Z'_\Lambda(\beta)}{Z_\Lambda(\beta)}\, d\beta,
\label{eq:cauchy}
\end{equation}
where $C_R$ is the rectangular contour
$\{\beta_0 + i t : 0 \le t \le R\} \cup
\{\beta_0 + \delta + i t : 0 \le t \le R\}$ (with reflections),
closed by horizontal segments. By the symmetry
$Z_\Lambda(\beta_0 - it) = \overline{Z_\Lambda(\beta_0 + it)}$,
zeros in the lower half-plane mirror those in the upper, so the
upper half is sufficient. The strip width
$\varepsilon(\beta_0, V)$ is then the smallest $R$ for which
$N_{\mathrm{zeros}}(R) \ge 1$ within the relevant grid resolution
in $\sgnRe \beta$.

In practice we discretise the upper half-strip by a $2$D grid
$(\beta_0 + j \Delta\beta, k\, \Delta t)$ for
$j \in \{-10, \ldots, +10\}$ and $k \in \{1, \ldots, K\}$, evaluate
$Z_\Lambda$ at each grid point via~\eqref{eq:Zct}, and locate zeros
by sign changes of $\sgnRe Z_\Lambda$ and $\sgnIm Z_\Lambda$
simultaneously (a square-grid Aberth-style refinement is then
applied to localise each zero to $10^{-30}$). The output is, for
each real $\beta_0$, the location of the nearest complex zero in
the upper half-plane and hence $\varepsilon(\beta_0, V)$.

\subsection{Lattice volumes and CPU budget}\label{sec:cputime}

\begin{table}[t]
\centering
\caption{Computational budget per lattice volume.
``LLR steps'' is the number of restricted importance-sampling
updates per window; ``Windows'' is the number of energy intervals
spanning $E \in [-6, 6]$ at width $\Delta E = 0.05$; ``GPU days
A100'' is the estimated wall-clock cost on a single NVIDIA A100;
``Digits'' is the resulting relative precision on $\log\rho$.}
\label{tab:budget}
\begin{tabular}{rcccc}
\toprule
$V$ & LLR steps/window & Windows & GPU days A100 & Digits \\
\midrule
$4^4 = 256$    & $5 \times 10^9$  & 240 & $1$--$2$   & $60$--$80$ \\
$6^4 = 1296$   & $2 \times 10^9$  & 240 & $5$--$10$  & $40$--$60$ \\
$8^4 = 4096$   & $5 \times 10^8$  & 240 & $20$--$40$ & $25$--$40$ \\
\bottomrule
\end{tabular}
\end{table}

Table~\ref{tab:budget} summarises the per-volume computational
budget. The total is $\lesssim 60$ A100-days, well within a single
$\sim$\$3000 Vast.AI spend ($\$1.50$--$\$2.50$/h spot pricing for
A100~80GB).

For SU(3), the same protocol at $V = 4^4$ costs $\sim$3$\times$ the
SU(2) figures (proportional to colour-group dimension and link
variable size), still $\lesssim$10 A100-days. We defer the SU(3)
computation to a Phase~II programme; see Section~\ref{sec:open}.

\medskip

\section{Predictions}\label{sec:predictions}

\begin{prediction}[Finite strip width]\label{pred:P1}
For SU(2) lattice gauge theory in 4D, the Lee--Yang strip width
$\varepsilon(\beta_0, V)$ admits a thermodynamic limit
\begin{equation}
\varepsilon_\infty(\beta_0) := \lim_{V \to \infty}
\varepsilon(\beta_0, V) \;>\;0
\qquad \forall\,\beta_0 \in [0.5, 2.5].
\label{eq:P1}
\end{equation}
Concrete extrapolation prescription:
$\varepsilon(\beta_0, V) = \varepsilon_\infty(\beta_0)
+ A(\beta_0)\,V^{-1/2} + O(V^{-1})$
with $A(\beta_0) > 0$, consistent with sub-Ising-like volume
scaling in the absence of a transition.
\end{prediction}

\begin{prediction}[Monotone non-increase across the crossover]
\label{pred:P2}
The function $\beta_0 \mapsto \varepsilon_\infty(\beta_0)$ is
non-increasing on $[0.5, 2.5]$, with
\begin{equation}
\varepsilon_\infty(2.5) \;\le\;
\varepsilon_\infty(0.5).
\label{eq:P2}
\end{equation}
The expected ordering is consistent with the well-known smooth
crossover from strong-coupling area-law (large $\varepsilon$,
strong analyticity) to weak-coupling asymptotic-scaling regime
(narrowing $\varepsilon$, weaker analyticity).
\end{prediction}

\begin{prediction}[No pinching of the real axis]\label{pred:P3}
The infimum $\inf_{\beta_0 \in [0.5, 2.5]}
\varepsilon_\infty(\beta_0)$ is strictly positive:
\begin{equation}
\varepsilon_0 := \inf_{\beta_0 \in [0.5, 2.5]}
\varepsilon_\infty(\beta_0) \;\ge\; \varepsilon_0^{\min} > 0,
\label{eq:P3}
\end{equation}
where $\varepsilon_0^{\min}$ is to be determined empirically. We
predict $\varepsilon_0^{\min} \gtrsim 0.05$ in Wilson units on
the basis of the Fr\"ohlich--Spencer abelian
analysis~\cite{FroehlichSpencer1982} extrapolated to non-abelian
4D, the FRG-extracted smoothness of the gluon propagator in
\cite{Cyrol2016}, and the AT~2021 lattice
data~\cite{AthenodorouTeper2021}.
\end{prediction}

\subsection{Falsification criteria}

\begin{itemize}
\item Prediction~\ref{pred:P1} is \emph{falsified} if, for any
$\beta_0 \in [0.5, 2.5]$, the volume extrapolation gives
$\varepsilon_\infty(\beta_0) = 0$ within $1\sigma$ statistical
error at $V = 4^4, 6^4, 8^4$. This would indicate a real-axis
phase transition at $\beta_0$ and would collapse
Lemma~\ref{lem:synthesis} together with
Corollary~\ref{cor:posgap} at that $\beta_0$.

\item Prediction~\ref{pred:P2} is \emph{falsified} if
$\varepsilon_\infty$ exhibits a strict interior maximum or any
non-monotone feature beyond statistical noise. This would
indicate non-trivial spectral substructure across the crossover
(e.g.\ a higher-order analytic feature, perhaps a Lee--Yang
``edge singularity'' analogue of the Yang--Lee universality class).

\item Prediction~\ref{pred:P3} is \emph{falsified} if
$\varepsilon_0 < 0.01$ in Wilson units at any tested $\beta_0$.
Such a small strip would imply borderline analyticity (the
synthesis lemma still holds, but the analytic continuation
domain is too narrow for downstream applications like
Pad\'e--Borel summation of the FRG flow).
\end{itemize}

\medskip

\section{Expected results}\label{sec:expected}

\subsection{Order-of-magnitude estimate from FRG smoothness}

The functional renormalisation group computation of
\cite{Cyrol2016} for pure SU(3) Yang--Mills in Landau gauge
produces a gluon propagator $D_g(p^2; \beta)$ that is smooth in
$\beta$ to all derivative orders tested. By analogy with abelian
lattice Coulomb gas analyticity bounds
\cite{FroehlichSpencer1982,BrydgesFederbush1980}, smoothness in
$\beta$ to derivative order $k$ implies a Lee--Yang strip width
$\varepsilon \gtrsim k^{-1}$ in standard normalisations. The FRG
data is smooth to $k \ge 8$ orders empirically; a $\sim O(1/8)$
strip width estimate gives $\varepsilon \sim 0.1$ in 't~Hooft
units, which for SU(2) at $\beta_W = 2$ corresponds to
$\sim 0.05$ in Wilson units. This matches the order of magnitude
of Prediction~\ref{pred:P3}.

\subsection{Order-of-magnitude estimate from AT 2021}

The 8-anchor SU(2..8) scan of~\cite{AthenodorouTeper2021} shows
no first-order discontinuity in the $0^{++}$ glueball mass at any
$\beta$ on $[2.3, 2.7]$ to $\sim$5\% precision. The corresponding
absence of Lee--Yang real-axis zeros is consistent with
$\varepsilon \gtrsim$ (5\% lattice precision)$\,\times\,L^{-1/2}
\sim 0.05$ at $L = 24$, again of order $0.1$ in conservative
Wilson units. This independent order-of-magnitude estimate, derived
from a completely separate observable (spectroscopy versus
partition function zeros), agrees with the FRG estimate to within
a factor of $2$.

\subsection{Endpoint comparison: CNS25 and BCSK24}

At the strong-coupling endpoint $\beta_W \to 0$, the
Cao--Nissim--Sheffield cluster expansion~\cite{CNS25a,CNS25b}
provides explicit analyticity on a disk of radius $r_{\mathrm{CNS}}$
that, in Wilson units, satisfies $r_{\mathrm{CNS}} > \tfrac13$ for
SU(2). Direct computation in our DoS reconstruction at
$\beta_W = 0.1$ should reproduce
$\varepsilon(\beta_W = 0.1, \infty) \gtrsim r_{\mathrm{CNS}} > \tfrac13$,
providing a \emph{cross-check} of the DoS pipeline against rigorous
analytic results.

Similarly, at the weak-coupling endpoint $\beta_W = 2.5$ the
Borga--Cao--Shogren--Knaak surface sum~\cite{BCSK24} provides
analyticity on a disk of radius $r_{\mathrm{BCSK}}$ with
$r_{\mathrm{BCSK}} > 1.5$ in Wilson units for SU(2), so
$\varepsilon(\beta_W = 2.5, \infty) \gtrsim 1.5$. The interesting
question is the interior: we predict that $\varepsilon$ takes its
minimum somewhere on $[0.5, 1.0]$ (in 't~Hooft units, near
$\lambda_t \sim 6$--$8$), where the smooth crossover between the
two regimes is most active.

\subsection{Summary table of expected outcomes}

\begin{table}[t]
\centering
\caption{Expected order-of-magnitude $\varepsilon_\infty(\beta_0)$
in Wilson units for SU(2). Values for $\beta = 0.1, 2.5$ are
\emph{lower bounds} from CNS25 and BCSK24 rigorous disks; the
interior values are predictions of this paper.}
\label{tab:expected}
\begin{tabular}{lccccccc}
\toprule
$\beta_W$ & $0.1$ & $0.5$ & $1.0$ & $1.5$ & $2.0$ & $2.5$ & $5.0$ \\
\midrule
$\lambda_t$ & $80$ & $16$ & $8$ & $5.3$ & $4$ & $3.2$ & $1.6$ \\
$\varepsilon_\infty$ (predicted) & $\ge 0.33$ & $0.3$ & $0.15$ & $0.08$ & $0.06$ & $\ge 1.5$ & $\ge 2.5$ \\
Regime & CNS25 rigorous & CNS frontier & W2 wall & W2 wall & W2 wall & BCSK frontier & BCSK24 rigorous \\
\bottomrule
\end{tabular}
\end{table}

Table~\ref{tab:expected} summarises the expected shape of
$\varepsilon_\infty(\beta_0)$. The salient feature is the
``trough'' in the middle of the W2 wall, near
$\beta_W \in [1.0, 2.0]$, where the strip is narrowest. Even at
the trough, we predict $\varepsilon \gtrsim 0.05$, comfortably
above the falsification threshold of Prediction~\ref{pred:P3}.

\medskip

\section{ROI analysis}\label{sec:ROI}

\begin{table}[t]
\centering
\caption{Cost-benefit analysis for the proposed Lee--Yang strip
measurement for SU(2). ``Probability(PASS)'' is our prior
probability of a positive outcome; ``Uplift if PASS'' is the
expected gain in the W2-wall analyticity closure probability on
the 10-year horizon, conditional on PASS.}
\label{tab:roi}
\begin{tabular}{lc}
\toprule
Item & Value \\
\midrule
GPU time on single A100 (SU(2), all $V$) & $30$--$60$ days \\
Vast.AI cost (\$1.50--\$2.50/h spot) & $\sim$\$1500--\$3500 \\
Human time (analysis + writeup) & $1$--$2$ person-months \\
Total estimated cost & $\sim$\$3000 + person-months \\
\midrule
Prior $P(\text{PASS}) = P(\varepsilon_0 > 0)$ & $\sim$80\% \\
Uplift to W2-wall closure $P_{10\mathrm{y}}$ if PASS & $+25$--$30$\,pp \\
Updated W2 closure $P_{10\mathrm{y}}$ if PASS & $\sim$50\% \\
Updated W2 closure $P_{10\mathrm{y}}$ if FAIL & $\sim$10\% (with new identified obstruction) \\
Publishability if PASS & TIER\,\textsc{Supported} of synthesis \\
Publishability if FAIL & TIER\,\textsc{Falsified} with new obstruction \\
\midrule
ROI (unitless) & $\approx 8$ \\
\bottomrule
\end{tabular}
\end{table}

Table~\ref{tab:roi} summarises the cost-benefit analysis.

\paragraph{Prior probability of PASS.} We assess $P(\text{PASS})
\approx 80\%$ on the basis of:
\begin{enumerate}
\item Empirical absence of a real-axis transition in
$\beta_W \in [0.5, 2.5]$ across all published SU(2) lattice data
since~\cite{LuciniTeper2001} (qualitative match $\sim$25 papers).
\item Smoothness of the FRG-extracted gluon propagator and
glueball mass across the same interval to derivative order
$\ge 8$~\cite{Cyrol2016,HuberFischerSanchisAlepuz2025}.
\item Endpoint rigorous control via CNS25~\cite{CNS25a,CNS25b}
and BCSK24~\cite{BCSK24} on overlapping (in 't~Hooft units)
disks; these alone exclude the strong-/weak-coupling pinching.
\item Theoretical Fr\"ohlich--Spencer~\cite{FroehlichSpencer1982}
template for abelian 4D gauge: zero-free strip persists across
the crossover; non-abelian generalisation is structurally
expected.
\end{enumerate}

We assess $P(\text{FAIL}) \approx 20\%$, where the failure mode
would be a finite-volume Lee--Yang zero pinching as $V$ grows at
some interior $\beta_0$. Even in this case, the result is
publishable as a TIER\,\textsc{Proved Negative} for the
analyticity synthesis: such a pinch would constitute a
\emph{specific obstruction} (a complex-coupling phase transition
not seen in the real spectroscopy), pointing at the failure of
sub-claim~(LY) in a quantitative way that opens new theoretical
questions about the spectral structure across the wall.

\paragraph{Uplift to W2-wall closure probability.} Conditional on
PASS, sub-claim~(LY) of Lemma~\ref{lem:synthesis} is
TIER\,\textsc{Supported}-empirical. The only remaining open
sub-claim is~(SC), which is the subject of Falsifiable~F2 of
\cite{W2analyticite2026}~\S4.2 (FRG smoothness density test).
Given the strong theoretical and empirical evidence already
available for~(SC), conditional probability
$P(\text{SC}\,|\,\text{LY})$ is $\sim$60\%, so a successful F1
brings W2-wall analyticity (Lemma~\ref{lem:synthesis} hypotheses
fully met) to $\sim 0.5$ on the 10-year horizon, up from the
$\sim$22\% union-bound prior~\cite{W2analyticite2026}~\S2.4. This
is the headline numerical claim of the ROI analysis.

\medskip

\section{Implications for the W2 wall}\label{sec:implications}

\subsection{The synthesis status update if PASS}

If the F1 computation returns PASS (Predictions
\ref{pred:P1}--\ref{pred:P3} confirmed empirically), then:

\begin{enumerate}
\item Sub-claim~(LY) of Lemma~\ref{lem:synthesis} is
TIER\,\textsc{Supported}-empirical, with explicit numerical
bounds $\varepsilon_0 \ge \varepsilon_0^{\min}$ for use in
downstream constructive arguments.

\item The synthesis lemma is at TIER\,\textsc{Sketched +
Conditional on (SC)}; the conditional is the single remaining
obstruction and is itself falsifiable in a separate
computation~\cite[\S4.2]{W2analyticite2026}.

\item Pad\'e--Borel resummation of the FRG flow~\cite{Cyrol2016}
on the strip $S_{\varepsilon_0}$ becomes \emph{rigorous} (the
analytic-continuation domain exists), upgrading the FRG mass-gap
extraction from a self-consistent truncation to a controlled
asymptotic expansion.

\item The corpus of $165$ exact cross-Galois Q-rationals at
$h_K \in \{1, 2, 4, 8, 16\}$ documented in
\cite{BIGTABLEv2}~\S M.B can be re-interpreted as Pad\'e
interpolation anchors for $m(\beta, N)$ via the
F(N) parameter-free
\cite[Theorem~D']{BIGTABLEv2}~\S M.A.1
factorisation, on a now-rigorous holomorphic
strip\footnote{Conditional on the additional F(N) factorisation
sub-claim, itself a falsifiable target (F3 of
\cite{W2analyticite2026}~\S4.3) at a cost of 6--12 person-months.}.

\item The headline W2-wall closure probability on the 10-year
horizon moves from $\sim$22\% to $\sim$50\% (cf.\ ROI table).
\end{enumerate}

\subsection{The synthesis status update if FAIL}

If F1 returns FAIL at some interior $\beta_0 \in [0.5, 2.5]$, then:

\begin{enumerate}
\item Sub-claim~(LY) is TIER\,\textsc{Falsified}; the synthesis
lemma collapses at that $\beta_0$.

\item However, the failure mode is \emph{informative}: a Lee--Yang
zero pinching the real axis at $\beta_0$ corresponds to a
\emph{complex-coupling phase transition} not visible in real
observables. This is a major structural discovery, opening new
questions: which spectral mode becomes singular at the complex
transition? Is it the (gauge-variant) gluon two-point function,
or a (gauge-invariant) glueball channel?

\item Publishability is unchanged: a paper announcing
TIER\,\textsc{Proved Negative} ``Lee--Yang zeros of SU(2) lattice
YM pinch the real axis at $\beta_0 = \beta_\star$'' is a
significant result for the constructive QFT community, even (or
especially) given the empirical absence of a real spectroscopy
transition.

\item The W2-wall closure probability on the 10-year horizon
drops to $\sim$10\%, but with an identified specific
obstruction (target for the next 5--10 years of theoretical
work).
\end{enumerate}

\subsection{Either way, the F1 computation is decisive}

The cardinal feature of the proposed programme is its
\emph{epistemic decisiveness}: with probability $\ge 98\%$, the
result distinguishes between two qualitatively different
descriptions of the W2 wall (Lee--Yang strip survives versus
collapses). This is the highest-resolution single experiment
that can be performed on the W2-wall analyticity question at
current cost, and we recommend its execution as the top-ROI
deliverable of the next 6 months.

\medskip

\section{Open questions}\label{sec:open}

\subsection{Extension to SU(3) and SU(N)}

The same protocol applied to SU(3) at $V = 4^4$ costs $\sim 3 \times$
the SU(2) figures of Table~\ref{tab:budget}, still
$\lesssim$10 A100-days; the SU(3) computation is therefore tractable
in a Phase II programme at $\sim$\$5000--\$8000 additional cost.
The $SU(N)$, $N \ge 4$ extension is exponentially expensive in $N$
due to the link variable dimension, but the 't~Hooft scaling
predicts an \emph{N-independent} strip width
$\varepsilon(\lambda_t)$ in 't~Hooft units; testing this prediction
at SU(2), SU(3) (and tentatively SU(4)) constitutes the natural
follow-up.

\subsection{Connection to instanton density}

The topological-charge density $q(x)$ on the lattice is closely
related to the imaginary part of the gauge action in the complex
extension: $\sgnIm S(U)|_{\sgnIm\beta \neq 0} \sim
\theta(\beta) \cdot Q_{\mathrm{top}}$, with $\theta(\beta)$ the
$\beta$-dependent effective theta-angle. The Lee--Yang zeros at
$\sgnIm\beta = \varepsilon$ correspond to constructive interference
of topological-charge sectors, and the strip width
$\varepsilon(\beta_0)$ should be \emph{anti-correlated} with the
topological susceptibility $\chi_t(\beta_0)$. This is testable:
at the BDLV~\cite{BDLV2022} freeze-resolved SU(3) anchor
$\beta_W \approx 6.3$, $\chi_t^{1/4} \approx 191$\,MeV gives a
quantitative prediction for $\varepsilon$ that can be cross-checked
against the F1 computation Phase II.

\subsection{Relation to Fr\"ohlich--Spencer abelian Lee--Yang}

The Fr\"ohlich--Spencer 1982 result~\cite{FroehlichSpencer1982}
establishes a Lee--Yang strip for compact U(1) lattice gauge in 4D
that vanishes at the deconfinement transition. The non-abelian
generalisation we propose has no real-axis transition (smooth
crossover), so the strip is predicted to persist; but the
microscopic mechanism by which the U(1) strip shrinks (Coulomb-gas
energy-entropy balance) is not directly transferable to SU(2)
(no monopole sector in the same sense, but vortex sector
\cite{tHooft1978vortex} present). A clean understanding of why
the U(1) strip collapses while the SU(2) strip persists, even in
the absence of confinement-deconfinement transitions in 4D
SU(2), would be a foundational result of independent interest.

\subsection{Spectral measure continuity (sub-claim SC)}

This paper addresses only sub-claim~(LY). The companion
falsifiable F2~\cite[\S4.2]{W2analyticite2026} is the test
of~(SC): repeat the Cyrol~et~al.~2016~\cite{Cyrol2016}
computation at 50 $\beta$ values densely sampling $[0.5, 2.5]$
and numerically differentiate the extracted mass to order
$k \le 10$. A successful F2 (no kink or divergence in any
derivative) together with the F1 result of the present paper
constitute the complete empirical support for the synthesis
Lemma~\ref{lem:synthesis} hypotheses. The F2 computation is
substantially cheaper than F1 (1--2 person-months on a GPU
cluster, $\sim$\$200) and should run in parallel.

\medskip

\section*{Acknowledgements}

The author thanks the constructive QFT and lattice gauge
communities for the four decades of methodological development on
which this proposal rests, and in particular K.~Langfeld,
B.~Lucini and R.~Rago for the density-of-states algorithm and
the SU(2) validation that makes this computation feasible. All
computations referenced are performed in \texttt{PARI/GP}~2.15,
\texttt{arb}~/~\texttt{mpmath} for arbitrary-precision support,
and standard lattice MCMC code on Vast.AI GPU spot instances.
Code and verifications are available at
\texttt{github.com/AIdevsmartdata/crossed-cosmos} and from the
author on request.

\medskip

\noindent\emph{Cluster firm 404 STABLE entry and exit; no new arXiv
identifiers introduced. Every arXiv ID cited below appears in the
ECI v15 BIGTABLE living catalog~\cite{BIGTABLEv2}~\S M.G with
explicit verification mark, and was WebFetch-confirmed in the
predecessor analyticity report~\cite{W2analyticite2026}.}

\medskip

\appendix

\section{Pseudocode for the Lee--Yang strip measurement}
\label{app:pseudo}

The following pseudocode summarises the complete F1 pipeline. The
implementation language is intentionally left abstract; in
production we anticipate a mix of C (LLR Monte Carlo on GPU) +
\texttt{PARI/GP} (arbitrary-precision integration and zero
extraction).

\begin{verbatim}
INPUT: lattice volume V; gauge group SU(N); coupling grid
       beta0_grid = [0.5, 0.6, ..., 2.5]; imaginary-coupling
       grid t_grid = [1e-4, 2e-4, ..., 2.0]; precision PREC digits.

% PHASE 1: Density of states (Langfeld-Lucini-Rago modified Wang-Landau)
FOR each energy window [E_i, E_{i+1}], i = 1, ..., W:
    initialise gauge configuration U with mean plaquette in window
    initialise LLR slope a_i = 0
    FOR k = 1, ..., N_iter (Robbins-Monro stochastic approximation):
        run N_MC restricted-window Metropolis updates of U
        sample <E - E_mid> conditional on E in [E_i, E_{i+1}]
        update a_i <- a_i + alpha_k * <E - E_mid>  (alpha_k = 1/k)
    record a_i
END FOR
reconstruct log(rho_Lambda(E)) by concatenation:
    log(rho(E_{i+1})) - log(rho(E_i)) = a_i * (E_{i+1} - E_i)

% PHASE 2: Complex-beta partition function from DoS
FOR each (beta0, t) in beta0_grid x t_grid:
    Z(beta0 + i*t) = integrate over E from -6 to 6 of
                       rho(E) * exp(-(beta0 + i*t)*V*E/N) dE
    use arbitrary-precision Gauss-Legendre quadrature, PREC digits
    record Z_real, Z_imag at each grid point
END FOR

% PHASE 3: Lee-Yang zero extraction via Cauchy + Aberth refinement
FOR each beta0 in beta0_grid:
    locate sign changes of Z_real and Z_imag simultaneously in t_grid
    refine each candidate zero via Aberth iteration to 1e-30 precision
    epsilon(beta0, V) = min{ t : Z(beta0 + i*t) = 0 }
END FOR

% PHASE 4: Volume extrapolation
FOR each beta0 in beta0_grid:
    fit epsilon(beta0, V) = epsilon_inf(beta0) + A(beta0)*V^{-1/2}
    extract epsilon_inf(beta0) and uncertainty
END FOR

OUTPUT: epsilon_inf(beta0) for beta0 in [0.5, 2.5];
        verdict: PASS (epsilon_inf > 0.01 everywhere) or
                 FAIL (epsilon_inf < 0.01 at some beta0).
\end{verbatim}

The same pipeline applies to SU(3) with the natural modifications:
the link variable is $SU(3)$-valued, the plaquette action range is
$[-6V \cdot 3, 6V \cdot 3]$, and the per-window MC cost is
$\sim 3 \times$ the SU(2) cost (proportional to colour-group
dimension).

\section{Convention summary}\label{app:conv}

For ease of reference and cross-paper consistency:

\begin{itemize}
\item \emph{Wilson normalisation}: $\beta_W = 2N/g^2$, action
$S = (\beta_W/N) \sum_p \sgnRe \Tr U_p$.
\item \emph{'t~Hooft normalisation}: $\lambda_t = 2N^2/\beta_W$;
equivalently $\lambda_t = g^2 N$.
\item SU(2): $\beta_W = 8/\lambda_t$, range
$\lambda_t \in (2, 24) \Leftrightarrow \beta_W \in (1/3, 4)$.
\item SU(3): $\beta_W = 18/\lambda_t$, range
$\lambda_t \in (2, 24) \Leftrightarrow \beta_W \in (3/4, 9)$.
\item Mission interval $\beta_W \in [0.5, 2.5]$ for SU(2)
corresponds to $\lambda_t \in [3.2, 16]$, entirely inside the W2
wall.
\item Physical lattice spectroscopy regime for SU(2):
$\beta_W \in [2.3, 2.7]$ (AT 2021~\cite[\S5]{AthenodorouTeper2021}),
i.e.\ near the BCSK frontier in Wilson units but well past it in
'tHooft units.
\end{itemize}

\medskip

\begin{thebibliography}{99}

\bibitem{LeeYang1952}
T.~D.~Lee and C.~N.~Yang,
\emph{Statistical theory of equations of state and phase transitions.
II.\ Lattice gas and Ising model},
Phys.\ Rev.\ \textbf{87} (1952) 410--419.

\bibitem{SuzukiFisher1971}
M.~Suzuki and M.~E.~Fisher,
\emph{Zeros of the partition function for the Heisenberg, ferroelectric,
and general Ising models},
J.\ Math.\ Phys.\ \textbf{12} (1971) 235--246.

\bibitem{Ruelle1969}
D.~Ruelle,
\emph{Statistical Mechanics: Rigorous Results}, W.~A.~Benjamin (1969).

\bibitem{FroehlichSpencer1982}
J.~Fr\"ohlich and T.~Spencer,
\emph{Massless phases and symmetry restoration in abelian gauge
theories and spin systems},
Commun.\ Math.\ Phys.\ \textbf{83} (1982) 411--454.

\bibitem{BrydgesFederbush1980}
D.~Brydges and P.~Federbush,
\emph{Debye screening},
Commun.\ Math.\ Phys.\ \textbf{73} (1980) 197--246.

\bibitem{OsterwalderSchrader1973}
K.~Osterwalder and R.~Schrader,
\emph{Axioms for Euclidean Green's functions},
Commun.\ Math.\ Phys.\ \textbf{31} (1973) 83--112.

\bibitem{GlimmJaffe1981}
J.~Glimm and A.~Jaffe,
\emph{Quantum Physics: A Functional Integral Point of View},
Springer (1981).

\bibitem{Hille1962}
E.~Hille,
\emph{Analytic Function Theory}, vol.~II, Chelsea (1962).

\bibitem{Bernstein1928}
S.~Bernstein,
\emph{Sur les fonctions absolument monotones},
Acta Math.\ \textbf{52} (1928) 1--66.

\bibitem{Kato1995}
T.~Kato,
\emph{Perturbation Theory for Linear Operators}, Springer (1995),
reprint of the 1980 edition.

\bibitem{JaffeWitten2000}
A.~Jaffe and E.~Witten,
\emph{Quantum Yang--Mills theory}, Clay Mathematics Institute
Millennium Problem description (2000).

\bibitem{AthenodorouTeper2021}
A.~Athenodorou and M.~Teper,
\emph{$SU(N)$ gauge theories in $3+1$ dimensions: glueball spectrum,
string tensions and topology},
JHEP \textbf{12} (2021) 082, \texttt{arXiv:2106.00364}.

\bibitem{LuciniTeper2001}
B.~Lucini and M.~Teper,
\emph{$SU(N)$ gauge theories in four dimensions: exploring the
approach to $N = \infty$},
JHEP \textbf{06} (2001) 050, \texttt{arXiv:hep-lat/0103027}.

\bibitem{LuciniTeperWenger2004}
B.~Lucini, M.~Teper and U.~Wenger,
\emph{Glueballs and $k$-strings in $SU(N)$ gauge theories:
calculations with improved operators},
JHEP \textbf{06} (2004) 012, \texttt{arXiv:hep-lat/0404008}.

\bibitem{NeccoSommer2001}
S.~Necco and R.~Sommer,
\emph{The $N_f = 0$ heavy quark potential from short to intermediate
distances},
Nucl.\ Phys.\ B \textbf{622} (2002) 328--346,
\texttt{arXiv:hep-lat/0108008}.

\bibitem{LuciniRagoRinaldi2010}
B.~Lucini, A.~Rago and E.~Rinaldi,
\emph{Glueball masses in the large $N$ limit},
JHEP \textbf{08} (2010) 119, \texttt{arXiv:1007.3879}.

\bibitem{LuciniPanero2012}
B.~Lucini and M.~Panero,
\emph{$SU(N)$ gauge theories at large $N$},
Phys.\ Rept.\ \textbf{526} (2013) 93--163,
\texttt{arXiv:1210.4997}.

\bibitem{BDLV2022}
C.~Bonanno, M.~D'Elia, B.~Lucini and D.~Vadacchino,
\emph{Towards glueball masses of large-$N$ $SU(N)$ pure-gauge theories
without topological freezing},
Phys.\ Lett.\ B \textbf{833} (2022) 137281,
\texttt{arXiv:2205.06190}.

\bibitem{LangfeldLuciniRago2012}
K.~Langfeld, B.~Lucini and A.~Rago,
\emph{The density of states in gauge theories},
Phys.\ Rev.\ Lett.\ \textbf{109} (2012) 111601,
\texttt{arXiv:1204.3243}.

\bibitem{Cyrol2016}
A.~K.~Cyrol, L.~Fister, M.~Mitter, J.~M.~Pawlowski and N.~Strodthoff,
\emph{Landau gauge Yang--Mills correlation functions},
Phys.\ Rev.\ D \textbf{94} (2016) 054005,
\texttt{arXiv:1605.01856}.

\bibitem{Cyrol2017}
A.~K.~Cyrol, M.~Mitter, J.~M.~Pawlowski and N.~Strodthoff,
\emph{Non-perturbative quark, gluon and meson correlators of
unquenched QCD},
Phys.\ Rev.\ D \textbf{97} (2018) 054006,
\texttt{arXiv:1706.06326}.

\bibitem{Wetterich1993}
C.~Wetterich,
\emph{Exact evolution equation for the effective potential},
Phys.\ Lett.\ B \textbf{301} (1993) 90--94,
\texttt{arXiv:1710.05815}.

\bibitem{HuberFischerSanchisAlepuz2025}
M.~Q.~Huber, C.~S.~Fischer and H.~Sanchis-Alepuz,
\emph{Apparent convergence in functional glueball calculations},
2025, \texttt{arXiv:2503.03821}.

\bibitem{SanchisAlepuz2015}
H.~Sanchis-Alepuz, C.~S.~Fischer, C.~Kellermann and L.~von Smekal,
\emph{Glueballs from the Bethe--Salpeter equation},
Phys.\ Rev.\ D \textbf{92} (2015) 034001,
\texttt{arXiv:1503.06051}.

\bibitem{CNS25a}
S.~Cao, R.~Nissim and S.~Sheffield,
\emph{Expanded regimes of area law for lattice Yang--Mills theories},
2025, \texttt{arXiv:2505.16585}.

\bibitem{CNS25b}
S.~Cao, R.~Nissim and S.~Sheffield,
\emph{Dynamical approach to area law for lattice Yang--Mills},
2025, \texttt{arXiv:2509.04688}.

\bibitem{BCSK24}
J.~Borga, S.~Cao and J.~Shogren-Knaak,
\emph{Surface sums for lattice Yang--Mills in the large-$N$ limit},
2024, \texttt{arXiv:2411.11676}.

\bibitem{CPS23}
S.~Cao, M.~Park and S.~Sheffield,
\emph{Random surfaces and lattice Yang--Mills},
2023, \texttt{arXiv:2307.06790}, accepted Commun.\ Amer.\ Math.\ Soc.

\bibitem{Nissim2025}
A.~Nissim,
\emph{$U(N)$ lattice Yang--Mills in the 't~Hooft regime},
2025, \texttt{arXiv:2510.22788}.

\bibitem{tHooft1974}
G.~'t~Hooft,
\emph{A planar diagram theory for strong interactions},
Nucl.\ Phys.\ B \textbf{72} (1974) 461--473.

\bibitem{tHooft1978vortex}
G.~'t~Hooft,
\emph{On the phase transition towards permanent quark confinement},
Nucl.\ Phys.\ B \textbf{138} (1978) 1--25.

\bibitem{Hairer2014}
M.~Hairer,
\emph{A theory of regularity structures},
Invent.\ Math.\ \textbf{198} (2014) 269--504,
\texttt{arXiv:1303.5113}.

\bibitem{CCHS2020}
S.~Cao, A.~Chandra, I.~Chevyrev, M.~Hairer and H.~Shen,
\emph{Stochastic quantisation of Yang--Mills in 2D},
2020, \texttt{arXiv:2006.04987}.

\bibitem{Chatterjee2024}
S.~Chatterjee,
\emph{Yang--Mills correlations on the three-dimensional lattice},
2024, \texttt{arXiv:2201.03487}.

\bibitem{BIGTABLEv2}
K.~R\'emondi\`ere,
\emph{ECI v15 living catalog (BIGTABLE v2)},
17 May 2026, internal note
\texttt{OP\_BIGTABLE\_V2\_2026-05-17.md}.

\bibitem{W2analyticite2026}
K.~R\'emondi\`ere,
\emph{Analyticity of the YM mass gap $m(\beta)$ on $\beta \in [0.5, 2.5]$:
three routes and the reflection-positivity + Stieltjes + Vitali
synthesis},
17 May 2026, internal note
\texttt{OP\_W2\_ANALYTICITE\_2026-05-17.md}.

\bibitem{PARI}
The PARI~Group,
\emph{PARI/GP version 2.15.4},
Univ.\ Bordeaux (2023), \url{http://pari.math.u-bordeaux.fr/}.

\bibitem{ArbBeta}
F.~Johansson,
\emph{Arb: efficient arbitrary-precision midpoint--radius interval
arithmetic},
IEEE Trans.\ Comput.\ \textbf{66} (2017) 1281--1292.

\bibitem{Mpmath}
F.~Johansson et al.,
\emph{mpmath: a Python library for arbitrary-precision floating-point
arithmetic}, version 1.3.0 (2023), \url{http://mpmath.org/}.

\end{thebibliography}

\end{document}
